\paragraph{Proofs - Separoid Axioms for $i\sigma$- and $id$-Separation}\hfill%

In the following let $G=(J,V,E,L)$ be a CDMG and $A,B,C,D \ins J \cup V$
(not necessarily disjoint) subsets of nodes.\\
Since $i\sigma$-separation can be expressed as $id$-separation in an acyclification $G'$
of $G$ (see Definition~\ref{def:cdg_acyclification}) by Proposition~\ref{prp:sigma_sep_as_d_sep}
%theorem \ref{thm-acy-sigma-sep}:
\[ A \isPerp_G B \given C  \qquad\iff\qquad A \idPerp_{G'} B \given C, \]
we can w.l.o.g.\ in the proof assume that we only use $id$-separation.\\

Recall that we say that \emph{$A$ is $id$-separated from $B$ given $C$} in $G$, in symbols:
    \[ A \idPerp_G B \given C, \]
 if every walk from a node in $A$ to a node in $J \cup B$ (sic!) is $d$-blocked by $C$. \\
 Again a walk $\pi$ is \emph{$d$-blocked by $C$} if it either contains a non-collider 
 (i.e.\ either an end node, fork, left/right chain) in $C$ or a collider not in $C$.

 We abbreviate the ternary relations in the following: $\displaystyle \iPerp := \idPerp_G$.


\begin{Lem}[Extended Left Redundancy]
    \[D \ins A \implies D \iPerp  B \given A.\]
\begin{proof}
    If $\pi$ is a walk from a node $v$ in $D$ to a node $w$ in $J \cup B$ then its first end node is in $A$, so $\pi$ is $d$-blocked by $A$.
\end{proof}
\end{Lem}

\begin{Lem}[$J$-Restricted Right Redundancy]
    \label{d-sep-restricted-right-redundancy}
    \[A \iPerp  \emptyset \given C \cup J \qquad \text{always holds.}\]
\begin{proof}
    If $\pi$ is a walk from a node $v$ in $A$ to a node $w$ in $J$ then its last end node is in $C \cup J$, so $\pi$ is $d$-blocked by $C \cup J$.
\end{proof}
\end{Lem}


\begin{Lem}[$J$-Inverted Right Decomposition]
    \[A \iPerp  B \given C \implies A \iPerp  J \cup B \given C.\]
\begin{proof}
 If $\pi$ is a walk from a node $v$ in $A$ to a node $w$ in $J \cup J \cup B$ then $w \in J \cup B$. %or $w \in C$.
 If $w \in J \cup B$ then by assumption $\pi$ is $d$-blocked by $C$.
 %If $w \in C$ then $w$ is an end node of $\pi$ that lies in $C$. So also in this case $\pi$ is $d$-blocked by $C$.
\end{proof}
\end{Lem}

\begin{Lem}[Left Decomposition]
    \[ A \cup D \iPerp  B \given C \implies D \iPerp  B \given C.\]
\begin{proof}
If $\pi$ is a walk from a node $v$ in $D$ to a node $w$ in $J \cup B$, then 
the walk $\pi$ is also a walk from $A \cup D$ to $J \cup B$, which by assumption is $d$-blocked by $C$.
\end{proof}
\end{Lem}

\begin{Lem}[Right Decomposition]
    \[ A \iPerp  B \cup D\given C  \implies A \iPerp  D \given C.\]
\begin{proof}
If $\pi$ is a walk from a node $v$ in $A$ to a node $w$ in $J \cup D$, then %, because $J \cup D \ins J \cup B$, 
the walk  $\pi$ is also a walk from $A$ to $J \cup B \cup D$, which by assumption is $d$-blocked by $C$. 
\end{proof}
\end{Lem}

\begin{Lem}[Left Weak Union]
    \label{d-sep-left-weak-union}
    \[A \cup D \iPerp  B  \given C \implies A \iPerp  B \given D\cup C.\]
\begin{proof}
    Let us assume the contrary:  $A \niPerp  B \given D\cup C$. Then there exists a shortest $(D \cup C)$-$d$-open walk $\pi$ 
    from a node $v$ in $A$ to a node $w$ in $J \cup B$ in $G$ such that every collider of $\pi$ is in $D \cup C$.
    Then every non-collider of $\pi$ is not in $D \cup C$.\\
    If now $\pi$ does not contain any node from $D \sm C$ then every collider of $\pi$ is in $C$. 
    This implies that $\pi$ is $C$-$d$-open, 
    which contradicts the assumption: $A \cup D \iPerp  B  \given C$.\\
    So we can assume now that $\pi$ contains a node in $D \sm C$. 
    Then consider the shortest sub-walk $\tilde \pi$ in 
    $\pi$ from $w \in J \cup B$ to a node $u \in D \sm C$. This means that $\tilde \pi$ does not contain any collider in $D\sm C$, 
    so they are all in $C$. So $\tilde \pi$ is $C$-$d$-open walk from $A \cup D$ to $J \cup B$. 
    This contradicts the assumption: $A \cup D \iPerp  B  \given C$.
\end{proof}
\end{Lem}

\begin{Lem}[Right Weak Union] 
    \[A \iPerp  B \cup  D \given C \implies A \iPerp  B \given D \cup  C.\]
\begin{proof}
    Follow the same steps as in Left Weak Union (Lemma~\ref{d-sep-left-weak-union}), 
    but this time get a contradiction with:
    $A \iPerp  B \cup  D \given C$.
   % Lets assume the contrary:  $A \niPerp  B \given D\cup C$. Then there exists a shortest $(D \cup C)$-$d$-open walk $\pi$ 
   % from a node $v$ in $A$ to a node $w$ in $J \cup B$ in $G$.
   % Then every collider of $\pi$ is in $D \cup C$ and every non-collider of $\pi$ is not in $D \cup C$.\\
   % If now $\pi$ does not contain any node from $D \sm C$ then every collider of $\pi$ lies in $C$. 
   % This implies that $\pi$ is $C$-$d$-open, 
   % which contradicts the assumption: $A \iPerp  B \cup D  \given C$.\\
    Then again we can assume that $\pi$ contains a node in $D \sm C$. 
    Then consider the shortest sub-walk $\tilde \pi$ in $\pi$ from $v \in A$ to a node $u \in D \sm C$.
    This means that $\tilde \pi$ does not contain any collider in $D\sm C$, 
    so they are all in $C$. So $\tilde \pi$ is $C$-$d$-open walk from $A$ to $J \cup B \cup D$. 
    This contradicts the assumption: $A \iPerp  B \cup D  \given C$.
\end{proof}
\end{Lem}

\begin{Lem}[Left Contraction] 
        \label{d-sep-left-contraction}
    \[ (A \iPerp  B \given D \cup  C) \land (D \iPerp  B \given C) \implies A \cup  D \iPerp  B \given C.\]
\begin{proof}
    Let us assume the contrary:  $A \cup D \niPerp  B \given C$.
    Then there exists a shortest $C$-$d$-open walk $\pi$ 
    from a node $v$ in $A \cup D$ to a node $w$ in $J \cup B$ in $G$ such that every collider of $\pi$ is in $C$. 
    So every non-collider is not in $C$ and $v$ is 
    the only node of $\pi$ that is in $(A \cup D) \sm C$ (otherwise $\pi$ could be shortened). \\
    Also $v$ cannot lie in $D\sm C$ as it would contradict the assumption: $D \iPerp  B \given C$.
    Thus $v \in A \sm C$ and $\pi$ is a walk from $A$ to $J \cup B$ whose colliders are in $C \ins D \cup C$ 
    and all non-colliders are not in $D \cup C$.
    %Furthermore, $\pi$ cannot have any non-colliders in $D$ as $\pi$ was assumed to be a shortest walk from $A \cup D$ 
    %to $J \cup B$. So all non-colliders of $\pi$ are outside of $C \cup D$. So $\pi$ must be $(D \cup C)$-$d$-open, which 
    But this contradicts the other assumption: $A \iPerp  B \given D \cup  C$.
\end{proof}
\end{Lem}

\begin{Lem}[Right Contraction]
  \label{d-sep-right-contraction}
    \[ (A \iPerp  B \given D \cup  C) \land (A \iPerp  D \given C)  \implies A \iPerp  B \cup  D \given C.\]
\begin{proof}
    Let us assume the contrary:  $A  \niPerp  B \cup D \given C$.
    Then there exists a shortest $C$-$d$-open walk $\pi$ 
    from a node $v$ in $A$ to a node $w$ in $J \cup B \cup D$ in $G$ such that  every collider of $\pi$ is in $C$.
    So every non-collider is not in $C$ and $w$ is 
    the only node of $\pi$ that is in $(J \cup B \cup D)\sm C$ (otherwise $\pi$ could be shortened). \\
    Also $w$ cannot lie in $D\sm C$ as it would contradict the assumption: $A \iPerp  D \given C$.
    Thus $w \in (J \cup B) \sm C$ and $\pi$ is a walk from $A$ to $J \cup B$ whose colliders all are in $C \ins D \cup C$ 
    and all non-colliders are not in $D \cup C$.
    But this contradicts the other assumption: $A \iPerp  B \given D \cup  C$.
\end{proof}
\end{Lem}

\begin{Lem}[Right Cross Contraction]
    \[ (A \iPerp  B \given D \cup  C) \land (D \iPerp  A \given C) \implies A  \iPerp  B \cup  D \given C.\]
\begin{proof}
    Verbatim the same as Right Contraction (Lemma~\ref{d-sep-right-contraction}), 
    only the first contradiction is with: $D \iPerp A \given C$.
\end{proof}
\end{Lem}

\begin{Lem}[Flipped Left Cross Contraction]
    \label{d-sep-flc-contraction}
    \[ (A \iPerp  B \given D \cup  C) \land (B \iPerp  D \given C) \implies B \iPerp  A \cup  D \given C.\]
\begin{proof}
    Let us assume the contrary:  $B  \niPerp  A \cup D \given C$.
    Then there exists a shortest $C$-$d$-open walk $\pi$ 
    from a node $v$ in $B$ to a node $w$ in $J \cup A \cup D$ in $G$ such that every collider of $\pi$ is in $C$.
    So every non-collider is not in $C$ and $w$ is 
    the only node of $\pi$ that is in $(J \cup A \cup D)\sm C$ (otherwise $\pi$ could be shortened). \\
    Also $w$ cannot lie in $(J \cup D) \sm C$ as it would contradict the assumption: $B \iPerp  D \given C$.
    Thus $w \in A \sm C$ and the walk $\pi$ (in reverse direction) is a walk from $A$ to $B$ whose colliders are all in 
    $C \ins D \cup C$ 
    and all non-colliders are not in $D \cup C$.
    But this contradicts the other assumption: $A \iPerp  B \given D \cup  C$.
\end{proof}
\end{Lem}


\begin{Lem}[$J$-Restricted Symmetry]
    %If $J=\emptyset$ then we have:
    \[A \iPerp  B \given C \cup J \implies B \iPerp A \given C \cup J.\]
\begin{proof}
  This follows from Flipped Left Cross Contraction (Lemma~\ref{d-sep-flc-contraction}) with $D=\emptyset$ and $C \cup J$ in place of $C$ together with
    $J$-Restricted Right Redundancy (Lemma~\ref{d-sep-restricted-right-redundancy}).
%If $J = \emptyset$ then every walk from $B$ to $J \cup A =A$ gives a walk from $A$ to $B=J \cup B$ in reverse direction. 
%Since being $d$-blocked by $C$ is not dependend on the direction of the walk the claim follows.
\end{proof}
\end{Lem}


\begin{Lem}[Left Composition]
    \[ (A \iPerp  B \given  C) \land (D \iPerp  B \given C) \implies A \cup D \iPerp  B \given C.\]
\begin{proof}
    Let $\pi$ be a walk from a node $v$ in $A \cup D$ to a node $w$ in $J \cup B$.
    If $v \in A$ then $\pi$ is $d$-blocked by $C$ by assumption: $A \iPerp  B \given  C$.
    If $v \in D$ then $\pi$ is $d$-blocked by $C$ by assumption: $D \iPerp  B \given  C$.
\end{proof}
\end{Lem}

\begin{Lem}[Right Composition]
    \[ (A \iPerp  B \given  C) \land (A \iPerp  D \given C) \implies A \iPerp  B \cup  D \given C.\]
\begin{proof}
    Let $\pi$ be a walk from a node $v$ in $A$ to a node $w$ in  $J \cup B \cup D$.
    If $w \in J \cup B$ then $\pi$ is $d$-blocked by $C$ by assumption: $A \iPerp  B \given  C$.
    If $w \in J \cup D$ then $\pi$ is $d$-blocked by $C$ by assumption: $A \iPerp  D \given  C$.
\end{proof}
\end{Lem}

\begin{Lem}[Left Intersection] Assume that $A \cap D = \emptyset$, then:
    \[ (A \iPerp  B \given D \cup  C) \land (D \iPerp  B \given A \cup C) \implies A \cup D \iPerp  B \given   C.\]
\begin{proof}
    Let us assume the contrary:  $A \cup D \niPerp  B \given C$.
    Then there exists a shortest $C$-$d$-open walk $\pi$ 
    from a node $v$ in $A \cup D$ to a node $w$ in $J \cup B $ in $G$ such that every collider of $\pi$ is in $C$.
    So every non-collider is not in $C$ and $v$ is 
    the only node of $\pi$ that is in $(A \cup D)\sm C$ (otherwise $\pi$ could be shortened). \\
    If $v \in A$ then by the disjointness of $A$ and $D$ we have that $v \notin D$. 
    Then $\pi$ is a walk from $A$ to $J \cup B$ whose colliders are in $C \ins D \cup C$ 
    and all non-colliders are not in $(D \sm C) \cup C =D \cup C$. This contradicts the assumption: $A \iPerp  B \given D \cup  C$.\\
    If $v \in D$ then similarly we get a contradiction: $D \iPerp  B \given A \cup  C$
\end{proof}
\end{Lem}

\begin{Lem}[Right Intersection] Assume that $B \cap D = \emptyset$, then:
    \[ (A \iPerp  B \given D \cup  C) \land (A \iPerp  D \given B \cup C) \implies A  \iPerp  B \cup D \given   C.\]
\begin{proof}
    Let us assume the contrary:  $A \niPerp  B \cup D \given C$.
    Then there exists a shortest $C$-$d$-open walk $\pi$ 
    from a node $v$ in $A$ to a node $w$ in $J \cup B \cup D $ in $G$ such that every collider of $\pi$ is in $C$.
    So every non-collider is not in $C$ and $w$ is 
    the only node of $\pi$ that is in $(J \cup B \cup D) \sm C$ (otherwise $\pi$ could be shortened). \\
    If $w \notin B$ then $w \in J \cup D$. In this case $\pi$ is a walk from $A$ to $J \cup D$ where every collider is in
    $C \ins B \cup C$ and all non-colliders are not in $(B\sm C) \cup C = B \cup C$. 
    So $\pi$ is a $(B \cup C)$-$d$-open walk from $A$ to $J \cup D$.
    This contradicts the assumption: $A \iPerp  D \given B \cup C$.\\
    If $w \notin D$ then $w \in J \cup B$. In this case $\pi$ is a walk from $A$ to $J \cup B$ where every collider is in
    $C \ins D \cup C$ and all non-colliders are not in $D \cup C$. 
    So $\pi$ is a $(D \cup C)$-$d$-open walk from $A$ to $J \cup B$.
    This contradicts the assumption: $A \iPerp  B \given D \cup  C$.\\
    Since $B \cap D = \emptyset$ there are no other cases ($B^\cmpl \cup D^\cmpl = J \cup V$) and we are done.
\end{proof}
\end{Lem}
